\chapter{Model selection and honest evaluation}
\label{ch:model-selection}

\begin{learningobjectives}
After this chapter, you should be able to:
\begin{itemize}
  \item state what an evaluation design estimates;
  \item use cross-validation without crossing preprocessing or dependence boundaries;
  \item separate model selection from final performance estimation; and
  \item make a promotion decision with uncertainty and error analysis.
\end{itemize}
\end{learningobjectives}

\begin{chapterconnection}
Part II has introduced several plausible learners. Choosing among them is itself
an adaptive statistical procedure. This chapter treats the evaluation protocol
as an estimator and the promotion decision as a versioned result, not as a
contest for the largest displayed score.
\end{chapterconnection}

\section{Name the estimand before the splitter}

An evaluation score estimates performance under a hypothetical repetition. In
ordinary $K$-fold cross-validation, rows are partitioned into folds; each fold
is predicted by a model trained on the others. The resulting estimate concerns
training sets somewhat smaller than the full sample and a population represented
by the fold construction.

Grouped folds estimate transfer to held-out groups. Forward or rolling temporal
splits estimate transfer from past training windows to later periods. Spatial or
domain splits can estimate transfer across sites. A convenient splitter cannot
repair an unclear deployment claim.

\section{A pipeline owns every fitted transformation}

Within each fold, fit imputation, scaling, category handling, feature selection,
dimensionality reduction, calibration, and the predictor using only the fold's
training portion. Apply the fitted sequence to its validation portion. A
pipeline makes this boundary executable and ensures that the same transformations
travel with the final artifact.

Static operations that use no sample information can be defined outside, but
their schemas and versions still belong to the contract. The test is not
whether an operation is called ``preprocessing''; it is whether it learns from
observations.

\section{Selection creates optimism}

Suppose many noisy performance estimates are compared and the largest is
reported. Even when every estimate is individually unbiased, the maximum tends
to be optimistic. Hyperparameter search, feature choice, threshold selection,
and repeated manual notebook experiments all consume validation information.

Nested evaluation separates these roles. Inner folds choose the candidate;
outer folds estimate the entire selection procedure on observations hidden from
that choice. A final locked test set can provide a single additional estimate
after development stops. It must not become a repeatedly queried leaderboard.

\conceptcheck{If the test set influences one feature choice, what has become the
new test set, and what evidence has been lost?}

\section{Search should express prior structure and budget}

A grid is useful for a small structured set. Random search often covers
high-dimensional spaces more efficiently when only some dimensions matter.
Adaptive search uses prior trials to choose later trials and must be included in
the selection boundary.

Search ranges should be expressed on natural scales, such as logarithmic ranges
for regularization strengths. Invalid combinations should be excluded by
construction. Record every trial, seed, split identifier, duration, failure,
metric, and artifact reference. A failed run is data about the process, not a
row to delete silently.

\begin{professionalbox}
Specify a compute budget and stopping rule before search. Reproducibility means
another run can reconstruct data, folds, code, configuration, and selection
logic; it does not require pretending that parallel floating-point execution is
bitwise identical on every device.
\end{professionalbox}

\section{Uncertainty and dependence}

Fold scores are not independent replicates because training sets overlap.
Reporting their standard deviation as if it were a simple standard error can be
misleading. Bootstrap methods must resample at the independent unit: cell,
patient, site, or time block rather than row when dependence lives at that level.

Uncertainty should match the decision. A confidence interval for mean loss does
not bound every subgroup or future shift. Report point estimates alongside
sample counts, variation by relevant group or period, and sensitivity to design
choices. Practical equivalence may matter more than a tiny nominal difference.

\section{Error analysis is structured, not anecdotal}

After obtaining held-out predictions, stratify errors by target range, feature
support, group, time, missingness, and operational context. Inspect examples to
form hypotheses, then test those hypotheses on data not used to invent them when
possible. Do not remove inconvenient cases without a protocol.

Learning curves compare performance as training size grows. A persistent gap
between training and validation behavior suggests variance or shift; uniformly
poor behavior suggests limited signal, representation, or misspecification.
These are diagnostic patterns, not universal proofs.

\begin{misconceptionbox}
\textbf{Misconception:} cross-validation eliminates overfitting. 
\textbf{Correction:} it estimates a procedure under its split assumptions.
Repeatedly adapting to its results can overfit the validation process, and an
inappropriate fold structure can estimate the wrong use setting precisely.
\end{misconceptionbox}

\begin{workedexample}{a promotion decision, not a leaderboard}
Four candidate families are evaluated for capacity regression. Inner grouped
temporal folds choose transformations and hyperparameters. Outer future-batch
folds estimate the complete choice procedure. The ensemble has the lowest mean
error, but the interval overlaps the ridge model and its tail latency exceeds
the batch-service budget.

The evaluation report promotes ridge because it beats the baseline by the
prespecified practical margin, satisfies batch-level regression limits, and has
stable resource use. The ensemble remains a challenger. The report includes
split hashes, data and code versions, candidate search spaces, all failed trials,
and a signed decision rationale.
\end{workedexample}

\section{Exercises}
\begin{enumerate}
  \item Design splits for unseen cells, future batches, and a new laboratory.
  State the estimand for each.
  \item Draw nested cross-validation including preprocessing and threshold
  selection.
  \item Explain why trial deletion and repeated seed selection can bias a report.
  \item Write promotion criteria containing a baseline margin, uncertainty,
  subgroup limits, latency, artifact size, and rollback readiness.
\end{enumerate}

\begin{chapterreview}
Evaluation estimates a named future use under a split design. Pipelines preserve
fold boundaries for every learned transformation. Selection creates optimism,
so nested or locked evaluation must separate choosing from estimating. Search
has a budget and provenance; uncertainty respects independent units; and a
promotion gate combines predictive, subgroup, and operational evidence.
\end{chapterreview}

\begin{notebooklink}
Lecture 15 notebook 00 treats selection as estimation. Companion laboratories
implement nested pipelines, bounded search, error analysis, uncertainty, and a
versioned promotion report.
\end{notebooklink}
